\begin{abstract}
\aidraft{%
We build and validate a market regime-detection instrument --- a two-state statistical jump
model on return-only downside features --- to an unusual standard, and then map, under strict
preregistered one-look discipline, exactly where its information is and is not usable for
decisions.
As a measurement device the label is exceptional: \stabJM{} agreement under $\pm 2$-year
training-window perturbation (versus \stabHMM{} for an HMM ensemble), \switchesYr{} switches
per year, and live operation to the present.
As a decision input it is dominated at every use we tested by a reactive estimator that never
names the regime: binary exposure timing loses \feeJMVT{}~bps/yr to volatility targeting;
covariance conditioning loses \feeBreact{}~bps/yr to a plain EWMA covariance; the filter's own
confidence margin loses to an EWMA-volatility signal in \dgpCells{} synthetic environments; and
cross-asset defensive rotation adds no risk-adjusted value beyond exposure, its drawdown
protection matched by volatility targeting.
A single mechanism, shown in simulation before any real-data run, unifies the nulls: regime
value at daily horizons is an information-timing claim, and a causal detector's recognition lag
lands it precisely where even a perfectly informed oracle loses to an estimator that pays no
lag toll.
The instrument-versus-strategy distinction, usually a footnote, is the thesis: measurement
validity is not decision value, and the applied regime-switching literature's positive claims
live in the gap between them.%
}
\end{abstract}
